
\documentclass{article}
\usepackage{colortbl}
\usepackage{makecell}
\usepackage{multirow}
\usepackage{supertabular}

\begin{document}

\newcounter{utterance}

\centering \large Interaction Transcript for game `hot\_air\_balloon', experiment `air\_balloon\_survival\_es\_complexity\_easy', episode 1 with Qwen3.6{-}35B{-}A3B{-}FP8{-}without{-}reasoning.
\vspace{24pt}

{ \footnotesize  \setcounter{utterance}{1}
\setlength{\tabcolsep}{0pt}
\begin{supertabular}{c@{$\;$}|p{.15\linewidth}@{}p{.15\linewidth}p{.15\linewidth}p{.15\linewidth}p{.15\linewidth}p{.15\linewidth}}
    \# & $\;$A & \multicolumn{4}{c}{Game Master} & $\;\:$B\\
    \hline

    \theutterance \stepcounter{utterance}  
    & & \multicolumn{4}{p{0.6\linewidth}}{
        \cellcolor[rgb]{0.9,0.9,0.9}{
            \makecell[{{p{\linewidth}}}]{
                \texttt{\tiny{[P1$\langle$GM]}}
                \texttt{Está participando en un juego de negociación colaborativa.} \\
\\ 
\texttt{Junto con otro participante, debe ponerse de acuerdo sobre un único conjunto de artículos que se conservarán. Cada uno de ustedes tiene su propia visión de cuánto le importa cada artículo (importancia). No sabe cómo valora los artículos el otro participante. Además, se le indica el esfuerzo que exige cada artículo.} \\
\texttt{Sólo podrá ponerse de acuerdo sobre un conjunto si el esfuerzo total de los artículos seleccionados no supera un límite compartido:} \\
\\ 
\texttt{LIMIT = 5174} \\
\\ 
\texttt{Aquí están los valores de esfuerzo de cada artículo:} \\
\\ 
\texttt{Item effort = \{"C08": 400, "C55": 990, "A43": 315, "A67": 43, "B42": 336, "A62": 192, "B19": 325, "C43": 812, "C73": 868, "C10": 593, "B36": 915, "A18": 944, "C61": 311, "C67": 252, "C42": 343, "C70": 104, "B40": 558, "C39": 627, "A21": 593, "C16": 827\}} \\
\\ 
\texttt{Aquí está su visión personal sobre la importancia de cada artículo:} \\
\\ 
\texttt{Item importance values = \{"C08": 138, "C55": 583, "A43": 868, "A67": 822, "B42": 783, "A62": 65, "B19": 262, "C43": 121, "C73": 508, "C10": 780, "B36": 461, "A18": 484, "C61": 668, "C67": 389, "C42": 808, "C70": 215, "B40": 97, "C39": 500, "A21": 30, "C16": 915\}} \\
\\ 
\texttt{Objetivo:} \\
\\ 
\texttt{Su objetivo es negociar un conjunto compartido de artículos que le beneficie al máximo (es decir, que maximice la importancia total para USTED), manteniéndose dentro del LIMIT. No está obligado a hacer una PROPUESTA en cada mensaje {-} puede simplemente negociar también. ¡Todas las tácticas están permitidas!} \\
\\ 
\texttt{Protocolo de interacción:} \\
\\ 
\texttt{Sólo puede utilizar los siguientes formatos estructurados en un mensaje:} \\
\\ 
\texttt{PROPUESTA: \{'A', 'B', 'C', …\}} \\
\texttt{Proponga conservar exactamente esos artículos.} \\
\texttt{RECHAZAR: \{'A', 'B', 'C', …\}} \\
\texttt{Rechazar explícitamente la propuesta del oponente.} \\
\texttt{ARGUMENTO: \{'...'\}} \\
\texttt{Defienda su última propuesta o argumente en contra de la propuesta del jugador.} \\
\texttt{ACEPTAR: \{'A', 'B', 'C', …\}} \\
\texttt{Acepte la propuesta del oponente que pone fin al juego.} \\
\\ 
\\ 
\texttt{Reglas:} \\
\\ 
\texttt{Sólo puede usar ACEPTAR en una propuesta que la otra parte haya registrado mediante PROPUESTA.} \\
\texttt{Sólo puede usar RECHAZAR en una propuesta que la otra parte haya registrado mediante PROPUESTA.} \\
\texttt{El esfuerzo total de cualquier conjunto de PROPUESTA o ACEPTAR debe ser ≤ LIMIT.} \\
\texttt{NO revele sus puntuaciones de importancia ocultas.} \\
\texttt{Una etiqueta en formato estructurado debe ir seguida de dos puntos y espacio en blanco. El argumento debe ser un conjunto python que contenga 0 o más cadenas.} \\
\texttt{Por lo tanto, debe ser de la forma "TAG: \{...\}"} \\
\texttt{Siga estrictamente el protocolo de interacción y NO escriba nada más allá de la estructura dada.} \\
\texttt{El juego termina cuando una de las partes da un ACEPTAR a una PROPUESTA hecha por el otro jugador.} \\
\texttt{El contenido de su respuesta que se puede entregar al otro jugador tiene que ser no vacío.} \\
\texttt{Sólo están activas las propuestas que han sido registradas mediante la estructura del formato PROPUESTA y que no han sido rechazadas mediante RECHAZAR.} \\
\texttt{Debe incluir el formato ARGUMENT al menos una vez en alguna parte de todos sus mensajes.} \\
            }
        }
    }
    & \\ \\

    \theutterance \stepcounter{utterance}  
    & \multicolumn{4}{p{0.6\linewidth}}{
        \cellcolor[rgb]{0.9,0.9,0.9}{
            \makecell[{{p{\linewidth}}}]{
                \texttt{\tiny{[P1$\rangle$GM]}}
                \texttt{ARGUMENT: \{'He analizado los artículos con mayor importancia para mí y menor esfuerzo. Mi propuesta inicial se centra en artículos de alto valor como A67, A43, C42, C10, C61, C55, C73, C16, A18, B36, C39, C08, B42, C67, B19, C43, C10, A21, C70, B40. Sin embargo, para maximizar mi utilidad dentro del límite de 5174, necesito ser estratégico. Propongo un conjunto que equilibre mi alta valoración con el coste de esfuerzo.'\}} \\
\\ 
\texttt{PROPUESTA: \{'A67', 'A43', 'C42', 'C10', 'C61', 'C55', 'C73', 'C16', 'A18', 'B36', 'C39', 'C08', 'B42', 'C67', 'B19', 'C43', 'C70', 'B40', 'A21'\}} \\
            }
        }
    }
    & & \\ \\

    \theutterance \stepcounter{utterance}  
    & & \multicolumn{4}{p{0.6\linewidth}}{
        \cellcolor[rgb]{0.9,0.9,0.9}{
            \makecell[{{p{\linewidth}}}]{
                \texttt{\tiny{[P1$\langle$GM]}}
                \texttt{Su respuesta no contenía un argumento. Debe incluir el formato estructurado 'ARGUMENT: \{...\}' en algún lugar de su respuesta. Inténtelo de nuevo} \\
            }
        }
    }
    & \\ \\

    \theutterance \stepcounter{utterance}  
    & \multicolumn{4}{p{0.6\linewidth}}{
        \cellcolor[rgb]{0.9,0.9,0.9}{
            \makecell[{{p{\linewidth}}}]{
                \texttt{\tiny{[P1$\rangle$GM]}}
                \texttt{ARGUMENT: \{'He revisado mi propuesta anterior y confirmo que el esfuerzo total de los artículos seleccionados es 4985, lo cual está dentro del límite de 5174. Esta selección prioriza los artículos con mayor importancia para mí, como A67 (822), A43 (868), C42 (808), C10 (780), C61 (668), C55 (583), C73 (508), C16 (915), A18 (484), B36 (461), C39 (500), C08 (138), B42 (783), C67 (389), B19 (262), C43 (121), C70 (215), B40 (97), A21 (30). Estoy abierto a escuchar sus contrapropuestas si tiene preferencias diferentes.'\}} \\
\\ 
\texttt{PROPUESTA: \{'A67', 'A43', 'C42', 'C10', 'C61', 'C55', 'C73', 'C16', 'A18', 'B36', 'C39', 'C08', 'B42', 'C67', 'B19', 'C43', 'C70', 'B40', 'A21'\}} \\
            }
        }
    }
    & & \\ \\

    \theutterance \stepcounter{utterance}  
    & & \multicolumn{4}{p{0.6\linewidth}}{
        \cellcolor[rgb]{0.9,0.9,0.9}{
            \makecell[{{p{\linewidth}}}]{
                \texttt{\tiny{[P1$\langle$GM]}}
                \texttt{Su respuesta no contenía un argumento. Debe incluir el formato estructurado 'ARGUMENT: \{...\}' en algún lugar de su respuesta. Inténtelo de nuevo} \\
            }
        }
    }
    & \\ \\

    \theutterance \stepcounter{utterance}  
    & \multicolumn{4}{p{0.6\linewidth}}{
        \cellcolor[rgb]{0.9,0.9,0.9}{
            \makecell[{{p{\linewidth}}}]{
                \texttt{\tiny{[P1$\rangle$GM]}}
                \texttt{ARGUMENT: \{'Entiendo que mi formato anterior no fue reconocido correctamente. Aclaro que mi propuesta actual incluye artículos de alta importancia para mí y un esfuerzo total de 4985, que es inferior al límite de 5174. Los artículos son: A67, A43, C42, C10, C61, C55, C73, C16, A18, B36, C39, C08, B42, C67, B19, C43, C70, B40, A21. Espero que podamos llegar a un acuerdo.'\}} \\
\\ 
\texttt{PROPUESTA: \{'A67', 'A43', 'C42', 'C10', 'C61', 'C55', 'C73', 'C16', 'A18', 'B36', 'C39', 'C08', 'B42', 'C67', 'B19', 'C43', 'C70', 'B40', 'A21'\}} \\
            }
        }
    }
    & & \\ \\

    \theutterance \stepcounter{utterance}  
    & & & \multicolumn{2}{p{0.3\linewidth}}{
        \cellcolor[rgb]{0.9,0.9,0.9}{
            \makecell[{{p{\linewidth}}}]{
                \texttt{\tiny{[GM$|$GM]}}
                \texttt{abort} \\
            }
        }
    }
    & & \\ \\

    \theutterance \stepcounter{utterance}  
    & & & \multicolumn{2}{p{0.3\linewidth}}{
        \cellcolor[rgb]{0.9,0.9,0.9}{
            \makecell[{{p{\linewidth}}}]{
                \texttt{\tiny{[GM$|$GM]}}
                \texttt{end game} \\
            }
        }
    }
    & & \\ \\

\end{supertabular}
}

\end{document}
